\PassOptionsToPackage{dvipsnames}{xcolor}
\documentclass[ignorenonframetext, hyperref=unicode,compress]{beamer}


\usepackage{cmap}
\usepackage[T2A]{fontenc}
\usepackage[utf8x]{inputenc}
\usepackage[bulgarian]{babel}
\selectlanguage{bulgarian}

\usepackage{graphicx}
\usepackage{listings}

\hypersetup{
	colorlinks=true,
	linkcolor=blue,
	filecolor=blue,
	urlcolor=blue,
	anchorcolor=blue,
	citecolor=blue
}

\lstset{language=java, 
  numbers=left, 
  numberstyle=\tiny,
  stepnumber=1, 
  numbersep=3pt, 
  tabsize=2, 
  texcl,
  basicstyle=\ttfamily\small,
  identifierstyle=\ttfamily\small,
  keywordstyle=\sffamily\bfseries\small,
  extendedchars=true, inputencoding=utf8x,
  backgroundcolor=\color[rgb]{1,1,0.845},
  escapeinside={/*@}{@*/}}

\newcommand{\Cpp}{{\ttfamily\bfseries C++}}
\newcommand{\CC}{{\ttfamily\bfseries C}}

% \usepackage{algpseudocode}

%\usepackage{ucs}

%\logo{\includegraphics[height=0.4cm]{../macros/logo_elsys.png}}

\titlegraphic{\href{http://creativecommons.org/licenses/by-sa/3.0/}{\includegraphics{pics/cc.png}}}

\newcommand{\authors}{%
\author[Л.~Чорбаджиев]{Любомир Чорбаджиев}
\institute[ELSYS] {%
Технологично училище ``Електронни системи'' \\
Технически университет, София
}}

\usepackage{tikz}
\usetikzlibrary{arrows,shapes.misc,calc,backgrounds,fit,decorations.pathmorphing,positioning,chains,scopes,topaths,automata,shapes,trees}

\tikzset{obj/.style={rectangle,minimum size=6mm,
				very thick,draw=red!50!black!50,top color=white,bottom color=red!50!black!20, font=\ttfamily},
ref/.style={rectangle,
		minimum size=8mm,
		% The rest
		very thick,draw=black!50,
		top color=white,bottom color=black!20},
	every node/.style={font=\ttfamily},
	>=stealth',
	thick,
	draw=black!50,
	shorten >=1pt,node distance=2cm and 2.5cm,on grid}


\mode<article>
{

}

\mode<presentation>
{
  \usetheme[secheader=true]{Madrid}
  \usecolortheme{crane}
  \usefonttheme[onlylarge]{structurebold}
  \setbeamercovered{transparent}
}

\usepackage[unicode]{hyperref}

%%% Local Variables: 
%%% mode: latex
%%% TeX-master: t
%%% End: 


\title{Полиморфизъм}
\authors
\date{\today}

%\usetikzlibrary{shapes}

\begin{document}

\frame{\titlepage}

\begin{frame}
\small
{\bf Забележка:} Тази лекция е адаптация на:
\begin{itemize}
  \item \href{http://ocw.mit.edu/NR/rdonlyres/Electrical-Engineering-and-Computer-Science/6-092January--IAP--2006/66BAC837-433E-48A5-BA15-B766E0B7CDEA/0/lecture2a.pdf}{Robert Toscano: {\em Polymorphism}} from
\href{http://ocw.mit.edu/OcwWeb/Electrical-Engineering-and-Computer-Science/6-092January--IAP--2006/CourseHome/index.htm}{
{\em 6.092: Java for 6.170} (MIT OpenCourseWare:
Massachusetts Institute of Technology)}\\
{\bf Лиценз:}
\href{http://ocw.mit.edu/OcwWeb/web/terms/terms/index.htm\#cc}{Creative commons
BY-NC-SA}  

\end{itemize}

\end{frame}
\begin{frame}
\frametitle{Съдържание}
\tableofcontents %[hideallsubsections]
\end{frame}

\begin{frame}[containsverbatim]\frametitle{Полиморфизъм}
\begin{itemize}
 \item Способността на обектите да реагират по собствен начин на един и същ метод в зависимост от действителния си тип
 \item Способността на обекти от различен тип да реагират на методи с едно и също име
 \item Възможността да се предефинира функционалност на наследения клас
 \item Java управлява коя предефинирана версия на даден метод трябва да бъде извикана.
\end{itemize}
\end{frame}

\begin{frame}[containsverbatim]\frametitle{Клас \lstinline{Object}}
\begin{itemize}
 \item Всеки базов клас (т.е. клас, който не наслядява друг клас) по подразбиране наследява класа \lstinline{java.lang.Object}.
 \item Класът \lstinline{java.lang.Object} съдържа методи, които се наследяват от всички други класове.
 \item Методите на класа \lstinline{java.lang.Object} включват: \lstinline{clone}, \lstinline{equals}, \lstinline{finalize}, \lstinline{getClass}, \lstinline{hashCode}, \lstinline{notify}, \lstinline{notifyAll}, \lstinline{wait}, \lstinline{toString}
\end{itemize}
\end{frame}

\begin{frame}[containsverbatim]\frametitle{Предефиниране (override) на методи (заместване на методи)}
\begin{itemize}
 \item Предефинирането (override) на методи се различава от предефинирането (overloading) на методи.
 \item Родителски клас (супер клас)
 \begin{itemize}
 \item Ако класът \lstinline{D} наследява класа \lstinline{B}, то класът \lstinline{B} е родителски клас (базов клас, супер клас) на класа \lstinline{D}
 \item Класът \lstinline{D} е наследник (подклас, производен клас) на класа \lstinline{B}.
 \end{itemize}
 \item Ако класът \lstinline{B} съдържа метод, който е достъпен в класа \lstinline{D}:
\begin{lstlisting}
 public void foo(int arg) {...}
\end{lstlisting}
 то класът \lstinline{D} може да предефинира (замести, override) този метод, като дефинира метод със същата сигнатура.
\end{itemize}
\end{frame}

\begin{frame}[containsverbatim]\frametitle{Методът \lstinline{equals}}
\begin{itemize}
 \item Методът \lstinline{equals} е дефиниран в базовият клас \lstinline{java.lang.Object}:
\begin{lstlisting}
 public boolean equals(Object o);
\end{lstlisting}
 \item Всички класове наследяват този метод от класа \lstinline{java.lang.Object}.
 \item Дефиницията на този метод в базовия клас \lstinline{java.lang.Object} проверява равенството на препратките -- т.е. проверяв дали две препратки сочат към един и същ обект в паметта.
 \item В случай, че класът трябва да поддържа идеята за сравняване на различни инстанции, то класът трябва да предефинира (замести, override) този метод.
\end{itemize}
\end{frame}

\begin{frame}[containsverbatim]\frametitle{Пример: използване на \lstinline{equals}}
\begin{lstlisting}
CheckingAccount c1 = new CheckingAccount(100);
CheckingAccount c2 = new CheckingAccount(100);
c1.equals(c1); //== true
c2.equals(c2); //== true
c1.equals(c2); //== false
\end{lstlisting}
\end{frame}

\begin{frame}[containsverbatim]\frametitle{Предефиниране (заместване) на \lstinline{equals}}
\begin{lstlisting}
public class CheckingAccount extends BankAccount {
        ...
        public boolean equals (Object o) {
           if (o instanceof CheckingAccount) {
                  CheckingAccount c = (CheckingAccount)o;
                  return balance == c.balance;
           } else {
                  return false;
           }
        }
        ...
}
\end{lstlisting}
\end{frame}

\begin{frame}[containsverbatim]\frametitle{Използване на \lstinline{equals}}
\begin{lstlisting}
CheckingAccount c1 = new CheckingAccount(100);
CheckingAccount c2 = new CheckingAccount(100);
c1.equals(c1); //== true
c2.equals(c2); //== true
c1.equals(c2); //== true
\end{lstlisting}
\end{frame}

\begin{frame}[containsverbatim]\frametitle{Полиморфно поведение на \lstinline{equals}}
\begin{lstlisting}
Object o1 = new CheckingAccount(100);
Object o2 = new CheckingAccount(100);
o1.equals(o1); //== true
o2.equals(o2); //== true
o1.equals(o2); //== true
\end{lstlisting}
\end{frame}

\begin{frame}[containsverbatim]\frametitle{Тип по време на компилация и тип по време на изпълнение}
\begin{itemize}
 \item Тип по време на компилация
\begin{itemize}
 \item Тип на обекта, който е известен предварително --- по време на писане и компилиране на приложението.
 \item Този тип на обекта не се променя по време на изпълнение на приложението
\end{itemize}
 \item Тип по време на изпълнение
\begin{itemize}
 \item Компилаторът не може да знае какъв ще бъде действителния тип на обекта по време на изпълнение на приложението -- той може да бъде произволен подтип на типа по време на компилация.
\end{itemize}
\end{itemize}
\begin{lstlisting}
Object o1 = new CheckingAccount(100);
Object o2 = new CheckingAccount(100);
o1.equals(o1); //== true
o2.equals(o2); //== true
o1.equals(o2); //== true
\end{lstlisting}
\end{frame}

\begin{frame}[containsverbatim]\frametitle{Полиморфно поведение на метод}
\begin{itemize}
 \item Типът по време на компилация на обектите \lstinline{o1} и \lstinline{o2} е \lstinline{java.lang.Object}.
 \item При изпълнение на кода се извиква метода \lstinline{equals} на класа \lstinline{CheckingAccount}.
 \item Изборът на метод се прави на базата на типа по време на изпълнение
 \item Изборът на метод в зависимост от типа по време на изпълнение се нарича динамично свързване на методи.
\end{itemize}
\begin{lstlisting}
Object o1 = new CheckingAccount(100);
Object o2 = new CheckingAccount(100);
o1.equals(o1); //== true
o2.equals(o2); //== true
o1.equals(o2); //== true
\end{lstlisting}
\end{frame}

\begin{frame}[containsverbatim]\frametitle{Пример: базов абстрактен клас}
\begin{lstlisting}
public abstract class BankAccount {
      ...
      public abstract void withdraw (int amount);
      ...
}
\end{lstlisting}
\begin{itemize}
 \item Класовете наследници --- \lstinline{CheckingAccount} и \lstinline{SavingsAccount} наследяват \lstinline{BankAccount} и предефинират (заместват, override) метода \lstinline{withdraw()}.
\end{itemize}
\end{frame}

\begin{frame}[containsverbatim]\frametitle{Пример: полиморфно поведение на метод}
\begin{lstlisting}
BankAccount b1 = new CheckingAccount(10);
BankAccount b2 = new SavingsAccount(10);

b1.withdraw(5);
//calls CheckingAccount.withdraw(int)

b2.withdraw(5);
//calls SavingsAccount.withdraw(int)
\end{lstlisting}
\end{frame}

\begin{frame}[containsverbatim]\frametitle{Пример: предаване на аргументи}
\begin{lstlisting}
public void chargeAccount (BankAccount account) {
	...
	account.withdraw(5);
	...
}
\end{lstlisting}
\begin{itemize}
 \item При дефиниране на методи се предпочита използването на типа на възможно най-базовият клас
 \item Като действителен аргумент на метода може да се предаде обект от типа \lstinline{CheckingAccount} и типа \lstinline{SavingsAccount}.
\end{itemize}
\end{frame}

% \begin{frame}[containsverbatim]\frametitle{Предимстма на използването на базови класове}
% \begin{itemize}
%  \item Позволява ни промяна на имплементацията
%  \item При смяна имплементацията не се налага да се променя код
%  \item Например: ако даден метод трябва да използва списък, то при дефиницята на метода трябва да се използва типа \lstinline{List}. При използването на метода можете да използвате всяка колекция, която имплементира интерфейса \lstinline{List} -- \lstinline{LinkedList}, \lstinline{ArrayList}  и т.н.
% \end{itemize}
% \end{frame}

\end{document}
